\documentclass[11pt]{article}
\usepackage[margin=1in]{geometry}
\usepackage{amsmath,amssymb,amsthm}
\usepackage{graphicx}
\usepackage{hyperref}

\newtheorem{theorem}{Theorem}
\newtheorem{corollary}{Corollary}
\newcommand{\C}{\mathbb C}

\title{Component propagation of the sun-dual resolvent inclusion}
\author{Partial answer to an open question in arXiv:2203.12765}
\date{11 August 2026}

\begin{document}
\maketitle

\noindent\textbf{Status.}
\texttt{candidate\_partial\_result\_likely\_valid\_needs\_human\_review}.

\begin{abstract}
Kruse and Schwenninger ask whether the restricted adjoint resolvent of a
bi-continuous semigroup always maps the sun dual into the domain of its
generator at every point of the original resolvent set. We prove that the set
of points where this inclusion holds is a union of connected components of
the resolvent set. It contains the component containing the right half-plane
beyond the growth bound. Thus the answer is affirmative whenever the
resolvent set is connected, and any counterexample must live in a different,
unseeded resolvent component.
\end{abstract}

\section{Source question and proof intuition}

Let \((X,\|\cdot\|,\tau)\) be a sequentially complete d-consistent Saks space,
let \((T(t))_{t\geq0}\) be a bi-continuous semigroup with generator
\((A,D(A))\), and use the source paper's sun dual \(X^\bullet\) and generator
\(A^\bullet\). For \(\lambda\in\rho(A)\), write
\[
 R(\lambda,A)^\bullet=R(\lambda,A)'|_{X^\bullet}
\]
whenever the restriction is viewed as an operator on the sun dual.

The second open question after Proposition 4.8 in \cite{source} asks whether
\begin{equation}\label{eq:target}
 R(\lambda,A)^\bullet X^\bullet\subset D(A^\bullet)
 \qquad(\lambda\in\rho(A))
\end{equation}
always holds. The source paragraph is reproduced in
Figure~\ref{fig:source}.

The key observation is local-to-global. At one point where
\eqref{eq:target} holds, Theorem 4.7(a) of the source identifies the restricted
adjoint resolvent with the genuine resolvent of \(A^\bullet\). A Neumann
series propagates this identity locally. Conversely, norm-resolvent
continuity and closedness of \(A^\bullet\) preserve the property under limits.
The good set is therefore relatively open and closed.

\section{The component theorem}

Define
\[
 \Omega=\bigl\{\lambda\in\rho(A):
 R(\lambda,A)^\bullet X^\bullet\subset D(A^\bullet)\bigr\}.
\]

\begin{theorem}[Component propagation]\label{thm:component}
The set \(\Omega\) is relatively open and relatively closed in \(\rho(A)\).
Consequently, \(\Omega\) is a union of connected components of \(\rho(A)\).
It contains the component containing
\(\{\lambda\in\C:\operatorname{Re}\lambda>\omega_0(T)\}\).
\end{theorem}

\begin{proof}
Fix \(\lambda_0\in\Omega\). By Theorem 4.7(a) of \cite{source},
\begin{equation}\label{eq:seed}
 \lambda_0\in\rho(A^\bullet),\qquad
 R(\lambda_0,A)'|_{X^\bullet}=R(\lambda_0,A^\bullet).
\end{equation}
Put \(R_0=R(\lambda_0,A)\) and
\(S_0=R(\lambda_0,A^\bullet)\). If
\(|\lambda-\lambda_0|\,\|R_0\|<1\), the Neumann formula gives
\[
 R(\lambda,A)
 =R_0\bigl[I+(\lambda-\lambda_0)R_0\bigr]^{-1}.
\]
Since \(R_0'\) leaves \(X^\bullet\) invariant and restricts there to
\(S_0\), taking adjoints and restricting the norm-convergent series yields
\[
 R(\lambda,A)'|_{X^\bullet}
 =S_0\bigl[I+(\lambda-\lambda_0)S_0\bigr]^{-1}
 =R(\lambda,A^\bullet).
\]
The final operator maps \(X^\bullet\) into \(D(A^\bullet)\), so a
neighbourhood of \(\lambda_0\) lies in \(\Omega\). Thus \(\Omega\) is
relatively open.

Now let \(\lambda_n\in\Omega\) and
\(\lambda_n\to\lambda\in\rho(A)\). For \(x^\bullet\in X^\bullet\), set
\[
 y_n=R(\lambda_n,A)'x^\bullet,
 \qquad y=R(\lambda,A)'x^\bullet.
\]
Norm-resolvent continuity gives \(y_n\to y\) in \(X'\). Each
\(y_n\in D(A^\bullet)\subset X^\bullet\), and \(X^\bullet\) is norm closed
in \(X'\), hence \(y\in X^\bullet\). Again by Theorem 4.7(a),
\[
 A^\bullet y_n=\lambda_n y_n-x^\bullet
 \longrightarrow \lambda y-x^\bullet
 \quad\hbox{in }X^\bullet.
\]
The operator \(A^\bullet\) is closed on \(X^\bullet\), so
\(y\in D(A^\bullet)\). This holds for every \(x^\bullet\), proving
\(\lambda\in\Omega\). Therefore \(\Omega\) is relatively closed.

A relatively clopen subset is a union of connected components. Proposition
4.8(a) of \cite{source} gives \eqref{eq:target} for every
\(\operatorname{Re}\lambda>\omega_0(T)\). That half-plane is connected, so
its entire component lies in \(\Omega\).
\end{proof}

\begin{corollary}[Connected resolvent set]
If \(\rho(A)\) is connected, then \eqref{eq:target} holds for every
\(\lambda\in\rho(A)\).
\end{corollary}

\begin{proof}
Theorem~\ref{thm:component} makes the nonempty set \(\Omega\) clopen in the
connected space \(\rho(A)\), so \(\Omega=\rho(A)\).
\end{proof}

\section{Upgrade attempts and remaining obstruction}

The theorem reduces a possible failure to a connected component of
\(\rho(A)\) disjoint from the right half-plane. Such a component may be a
bounded spectral hole, although boundedness is not asserted in general.
Four focused upgrades were pursued.

\begin{enumerate}
\item A Riesz--Runge functional-calculus argument propagates through the
seeded component but cannot cross spectrum separating it from another
component; the missing pole is precisely the obstruction.
\item Abstract invariant-subspace examples can make a restriction acquire
extra spectrum (the bilateral/unilateral shift mechanism), but no realization
was found in the special norming sun-dual subspace of a d-consistent Saks
space.
\item Standard shift, multiplication, and translation semigroups examined in
this setting have matching resolvent behaviour and do not produce a failure.
\item An ordinal endpoint-projection example initially appeared to answer the
separate continuity-necessity question. It is excluded by d-consistency: its
endpoint functional is sequentially mixed-continuous but not mixed-continuous,
while condition (ii) applied to the bounded weak-null sequence formed by its
scalar multiples would force mixed continuity.
\end{enumerate}

Thus the unresolved mechanism is extra spectrum of the restricted adjoint on
an unseeded resolvent component, under all the d-consistency and sun-dual
constraints. Further local resolvent manipulations do not supply a seed there.

\paragraph{Novelty and review.}
The four cheap run indexes and exact-phrase/keyword searches found no prior
answer to the open question or statement of Theorem~\ref{thm:component}.
Because the proof is short, it may nevertheless be unindexed folklore;
novelty confidence is moderate. Expert review should check the restriction of
the Neumann series and the closed-graph limit.

\paragraph{Limitations.}
This is not a full answer in the disconnected case. It also does not settle
whether mixed-topology continuity of \(R(\lambda,A)\) is necessary for the
domain inclusion.

\begingroup
\raggedright
\begin{thebibliography}{9}
\bibitem{source}
K. Kruse and F. L. Schwenninger,
\emph{Sun dual theory for bi-continuous semigroups},
arXiv:2203.12765 (2022; revised 2023).
\end{thebibliography}
\endgroup

\begin{figure}[ht]
\centering
\includegraphics[width=\textwidth]{figures/open_questions_crop.png}
\caption{Source page 18: the two open questions following Proposition 4.8.}
\label{fig:source}
\end{figure}

\end{document}
